% From mitthesis package
% Version: 1.04, 2023/10/19
% Documentation: https://ctan.org/pkg/mitthesis


\chapter{Introduction}
% title:

% Smarter Stays, Smaller Footprints: A Data-Driven Solution for Sustainable Tourism in the Hospitality Industry

% Smarter Stays, Smaller Footprint: reducing energy and water usage in the sharing economy hospitality business in Costa 
% Rica through a data-driven solution

% Data collection and visualization to reduce the usage of energy and water in the sharing economy hospitality business
% in Costa Rica

% keywords:
% sustainable tourism
% smart home technology
% energy water efficiency
% energy monitoring
% water usage
% footprint
% sustainable practices
% hospitality industry
% Airbnb
% IoT
% internet of things
% smart devices

% 1. Trends and patterns
% - theories
% - methods
% - results
% 2. Themes
% 3. Debates and conflicts
% 4. Pivotal publications
% 5. Gaps

% Structure:
% Chronological
% Thematic ****
% Methodological
% Theoretical

\section{Overview}
% Comments from Martyn

% Make sure to include a reference to Airbnb  when first mentioned (use their home site).

% Although an interesting anecdote and a good motivator for your project idea, this should be avoided in academic
% writing for a software development project. Consider rewriting and talking more generally, i.e. focus on just
% detailing the house specifically, you can say "The first house, was constructed  in Costa Rica in a tropical forest
% with sustainability in mind. This included, low energy appliances, energy efficient lighting ....." etc. In short if
% you remove the personal pronouns and write in the passive you can still focus on this motivation for the project, but 
% in a more objective and academic way.

The tourism industry is booming with 1.3 billion international visitors in 2019 \cite{UNWTO2024a} and reaching 89\% of
pre-pandeminc levels in the first quater of 2024 \cite{UNWTO2024b}. This has a direct impact on traveling carbon 
footprint, not just on how people get to their destinations but also from their accommodation
\cite{Becken2001, Scott2008}. Energy consumption from hotels and other types of accommodation is one of the most
intensives \cite{Jos2016}, forcasted to double by 2040 \cite{Gossling2015}.

Bed and Breakfast style accomodation, particularly peer-to-peer has had a strong growth driven by the popularity of
online platforms like AirBnB \cite{AirBnB2024a} with over 8 million listings globally \cite{AirBnB2024b}. Is in this
type of accommodation where guests have more autonomy over the use of the facilities, and as a result the amount of 
resources consumed. Nevertheless, less research has been done in this area compared to hotels, despite the potential
\cite{Becken2016, Warren2017}.

Previous research has shown that electricity accounts for about 70\% of the total energy consumption in tourist
accomodation \cite{Becken2001, Warren2017}. Water usage is also a concern, with a strong nexus between water and
energy which could be reduced through water conservation as it can be a limited resource in some locations 
\cite{Becken2017, Lam2016}.

Influencing guest behaviour during their stay can be a challenge since no one can really see how guests are consuming
resources \cite{Juvan2014}. 

To be able to capture small changes in resource use, utilities have to be monitored at a high frequency \cite{Jos2016}. 
This is where intelligent technologies for monitoring come in, giving both the owner and the guest 'real time' readings
at high resolutions \cite{Warren2018a}.

\section{Problem Definition (Presentation of the Problem)}
% Comments from Martyn

% Problem definition

% Part of your argument is too stretched. There may be no solutions to monitor resources of water and electricity for
% Airbnb guests specifically, but solutions like the Google Nest, Smart Electricity and water meters do exist here that
% connect to a mobile  application straightforwardly, and even monitor gas usage. It might not be the same where you are
% though, but recognise that these solutions do exist. Also recognise, that like these devices, you will not be making
% this system interoperable. These are off the shelf solutions that you have pulled together, and not sensor devices you
% have developed and provided open access to. So it is likely that your solution will also need to keep the data
% private, and will unlikely be interoperable with other solutions either unless you can demonstrate this concretely. An
% initiative to supply the technology is unlikely to forthcoming from Airbnb due to the scale of the issue. Would hosts
% be happy being forced into paying for these devices and setting up them up correctly. This is not a good argument to
% support the development of your app. You can say that generally work in this area has been very little, and so this is
% why your project has impact.

% Focus instead on the simpler problem of making energy usage more transparent through your mobile application for
% Airbnb guests staying at your specific locations. In essence you are creating a framework and approach to easily
% providing, those who are energy conscious, the tools they need to adjust their behaviour whilst being hosted.

Thus, a big concern is that guests are not aware of their (real time) resource use when staying at an self-catering
style accomodation \cite{Juvan2014}.

Research has shown that people can cut their shower time by 27\% when they are given feedback on their water usage
\cite{Stewart2013}.

Studies conducted showed a reduction in electricity consumption by 27.4\% and 21.6\% in water usage when guests were
aware of their consumption from previous days \cite{Warren2018b}.

Make energy and water usage more transparent for AirBnB guests to adjust their behavour while being hosted.

Befefit for hosts: 

It might help address the issue of guests sleeping poorly due to uncomfortable inside temperatures \cite{Pallesen2016}.

The focus of this project is on shared accomodation made available through online platforms like AirBnB.

\section{Background Research}
% Comments from Martyn

% Background research 

% While the literature review provides a good overview of relevant topics related to Airbnb, sustainable tourism, and
% smart home technology, it could benefit from more in-depth exploration of existing solutions in the field. Try to
% include more recent research studies, industry reports, and academic articles to support your project's context and
% relevance.

% For example, this journal paper has over 800 citations so it likely to be a good reference, you might also find
% further references related to your discussion in the literary review that you can include.
% https://www.emerald.com/insight/content/doi/10.1108/JTF-11-2015-0048/full/html

% You might also include this paper, which discusses how guests feel about the use of smart devices at their
% accommodation and respond to the points raised.
% https://dl.acm.org/doi/abs/10.1145/3334480.3382900

% This paper, might also be worth reading as it touches on many of the topics you have discussed such as the shared
% economy and the potential benefit of incorporating smart technologies within this context. They also mention that the
% use of smart technology can be beneficial to urban areas.
% https://www.mdpi.com/2075-5309/13/9/2182

% Provide a reference Sustainable Hosting Plan, and Energy Saving Trust next to their first mention.

% You need a reference where you mention the 1990's being the start of IoT systems being developed thanks to the
% internet (See: https://www.researchgate.net/publication/356876316_A_Noble_Proposal_for_Internet_of_Garbage_Bins_IoGB).

% Introduction

% - Restate your research question
% - Emphasize the timeliness of the topic
% e.g. "many recent studies have focused on the problem of X"
% - Highlight gap in the literature
% e.g. "while there has been much research on X, few researchers have taken Y into consideration"

% Main body
% 1. Summarize and synthesize
% 2. Analize and interpret
% 3. Critically evaluate
% 4. Structure paragraphs

\subsection{Shared Economy Systems}

Shared economy is defined as “a socioeconomic system that enables an intermediated set of exchanges of goods and
services between individuals/organizations which aim to increase efficiency and optimization of underutilized resources
in society” \cite{Munoz2017}.

\subsection{Technology Applications in Self-Catering Accomodation}

\subsection{Smart Monitoring and the Internet of Things}

Advanced Metering Infrastructure (AMI) is a system that allows two-way communication between a utility and a customer,
it helps consumers monitor their energy and water consumption in near real-time. They offer a range of benefits, from 
accurate consumer billing, less cost of operation and features like leakage detection \cite{Monks2019}.

\subsection{Research Need and Knowledge Gap}

The ecosystem of the shared economy must be adapted to make information more transparent to guests and hosts. This in 
turn will enhance trust among skateholders \cite{Deng2021}.

There is a need to develop intelligent technology platforms that allow provider and users of shared services to make
make trustworthy decisions \cite{Hamari2016}.

Local official statistics at district level on tourists staying in shared accomodation are very scarce \cite{Johns2007}.

In 2022 the United Kingdom's Office for National Statistics (ONS) regognised the need for more novel data sources from 
booking portals, mobile phones, payment cards and travel agencies or any other data that can contribute to the
production of tourism statistics \cite{ONS2022}.

\section{Aim and Objectives}
% Comments from Martyn

% Aims and Objectives

% Try to be more specific with your objectives, remember you will revisit these as part of the critical evaluation of
% the final report to show that you have fulfilled them. When you say create a robust system, talk more about the
% sensors and how they are connected. i.e. 

% - Create an array of sensors monitoring water, energy usage and harvest the data via a Arduino / raspberryPI / PC /
% mobile phone.

% - Process, and analyse the data to generate statistics of minimum, maximum, and average daily usage of water, energy,
% and the inside and outside temperature.

% (You can't complete the second objective since it can only be achieved if you can demonstrate with evidence that the
% way you process, analyse and compare data will reduce carbon footprint of Airbnb stays specifically -- you would need 
% a period of monitoring to confirm this over a sample of people staying, you also do not know whether they have reduced
% their footprint if you do not know their baseline usage of these resources either at their own home, or over multiple
% stays at yours). In short, if you can't support a claim try to find another way to phrase it or provide
% references/evidence.

% The other objectives are specific and actionable.

% The main aim of this project is to develop a data-driven solution that educates tourists about the carbon footprint of
% Airbnb stays, empowers hosts to optimize resource usage, and promotes sustainable practices in the hospitality
% industry.
